
\chapter{Battery Lessons and Domain Interpretation}
\label{ch:battery-lessons}

This chapter steps back from the numerical results to articulate what the
battery domain \emph{means} for the ORIUS programme.  The battery is not
merely a convenient test-bed: it is a \emph{witness domain}---a setting
rich enough to expose the phenomena that ORIUS addresses, yet constrained
enough to permit rigorous measurement.

\section{Battery as a Witness, Not a Toy}

A witness domain must satisfy three criteria:
\begin{enumerate}
  \item \textbf{Physical fidelity.}  The plant dynamics (SOC evolution,
    efficiency losses, power/ramp constraints) are governed by real
    thermodynamic and electrical laws, not synthetic rules.
  \item \textbf{Telemetry realism.}  The fault injection model
    (packet dropout, delay jitter, covariate drift, stale sensors)
    mirrors failures observed in deployed SCADA and IoT systems.
  \item \textbf{Safety materiality.}  Constraint violations have
    measurable economic and physical consequences: over-discharge
    accelerates cell degradation, and SOC excursions can trigger
    protective shutdowns that disrupt grid services.
\end{enumerate}

The CPSBench battery plant satisfies all three criteria.  It is not a
``toy'' in the pejorative sense: the dynamics are calibrated to a
100~MWh / 50~MW grid-scale BESS, and the fault magnitudes are drawn
from field telemetry audits of operating installations.

\section{What the Battery Proof Surface Establishes}

The battery evaluation provides a \emph{proof surface}: a manifold of
evidence linking theoretical claims to empirical measurements.  The
surface consists of:

\begin{description}
  \item[Existence evidence.]  The OASG phenomenon (observed-safe but
    truly-unsafe actions) exists and occurs at non-trivial rates
    (TSVR~$= 3.9$\% for the deterministic LP).
  \item[Safety evidence.]  The DC3S shield eliminates the OASG gap
    (TSVR~$= 0$\% across all tested conditions).
  \item[Bound evidence.]  The expected violation count satisfies the
    ORIUS core bound $\mathbb{E}[V] \le \alpha (1 - \bar{w})\,T$
    with empirical slack.
  \item[Necessity evidence.]  Ablation of reliability-aware components
    re-introduces violations, confirming the No-Free-Safety theorem.
  \item[Temporal evidence.]  Certificate half-life and graceful
    degradation extend the safety guarantee into the time dimension.
\end{description}

\section{What Battery Cannot Establish Alone}

The battery domain is one domain.  It cannot, by itself, establish:
\begin{itemize}
  \item \emph{Universality.}  The claim that ORIUS applies to arbitrary
    CPS domains requires either additional witness domains or a
    domain-transfer argument (Chapter~\ref{ch:temporal-extensions}).
  \item \emph{Optimality.}  The intervention rate of 2.8\% is an
    empirical outcome, not a proven lower bound on the cost of safety.
  \item \emph{Adversarial completeness.}  The threat model
    summarised in Chapter~\ref{ch:outside-current-evidence} is bounded; a sufficiently
    powerful adversary who controls the physical plant is out of scope.
  \item \emph{Long-horizon stationarity.}  The 168-hour episodes are
    long enough for statistical significance but short relative to
    seasonal and multi-year battery degradation cycles.
\end{itemize}

\section{Why Battery Is the Right Reference Domain}

Among candidate CPS domains (HVAC, water networks, autonomous vehicles,
industrial process control), battery energy storage offers a unique
combination of advantages:

\begin{enumerate}
  \item \textbf{Scalar safety predicate.}  The SOC constraint
    $\text{SOC}_{\min} \le x_t \le \text{SOC}_{\max}$ is a clean,
    one-dimensional safety condition that admits exact verification.
  \item \textbf{Linear dynamics.}  The SOC transition
    $x_{t+1} = x_t + \eta\,P\,\Delta t$ is affine, enabling closed-form
    derivations that would require approximation in nonlinear domains.
  \item \textbf{Economic materiality.}  Battery arbitrage is a real
    market with quantifiable revenue, providing a natural cost metric
    for the price of safety.
  \item \textbf{Telemetry vulnerability.}  BESS deployments are
    among the most telemetry-dependent grid assets, making the
    degraded-observation setting practically motivated.
  \item \textbf{Data availability.}  Public grid datasets (EIA-930,
    ENTSO-E) provide realistic load, price, and renewable signals.
\end{enumerate}

\section{How Battery Lifts into ORIUS}

The battery results instantiate the ORIUS framework at a specific
domain point.  The lifting argument proceeds as follows:

\begin{enumerate}
  \item Every theorem in Part~III is stated in domain-general notation
    (state $x_t$, action $a_t$, safety predicate $\phi$) and then
    \emph{instantiated} in battery terms (SOC, dispatch, SOC bounds).
  \item The empirical results in Part~II provide the battery instantiation
    data that satisfies the theorem preconditions.
  \item To transfer to a new domain, one must verify that the new
    domain satisfies the same assumptions (A1--A8) and provide
    analogous empirical instantiation.
\end{enumerate}

This structure ensures that the battery evidence is neither under-claimed
(it does prove the theorems in one domain) nor over-claimed (it does not
prove them in all domains without further work).

\section{Chapter Summary}

The battery domain is a principled witness for the ORIUS safety framework.
Its physical fidelity, telemetry realism, and economic materiality make it
the strongest single-domain evidence base available.  The proof surface it
provides---spanning existence, safety, bounds, necessity, and temporal
extension---anchors the theoretical programme while honestly
acknowledging the limits of single-domain evidence.
